%!TEX root = ../../main.tex
%% 1.5 연구 질문 — M-RQ1–4 질문 본문은 고정 문구 (thesis/tools/mrq_hashes.json으로 검사)
%% 2026-09-22 M-RQ3 개정: "학습기의 선택" 절을 제거 (2장의 범위 선언과 충돌). 저장소의 mrq_hashes.json 재생성 필요.
\section{연구 질문}
\label{sec:intro-rq}

네 연구 질문은 절차의 네 단계와 제 \ref{ch:data}--\ref{ch:validation} 장에 하나씩 대응한다.

\begin{description}[leftmargin=3.6em,style=nextline]
  \item[M-RQ1 (교전 사례 구성)] 서로 다른 시간 해상도의 공개 상태·사건 기록에서 킬 조건부 교전 사례를 어떤 근거로 구성할 수 있으며, 경계와 규모 기준의 선택은 사례 구성에 어떤 변화를 만드는가?
  \item[M-RQ2 (결과 가치 구성)] 최종 경기 승패와 연결한 가치모형을 이용하여 교전 전후 전략적 가치 개선을 어떻게 정의할 수 있으며, 이 측정은 관측된 결과와 얼마나 대응하고 모형·구간 선택에 얼마나 의존하는가?
  \item[M-RQ3 (사전 정보와 예측)] 이렇게 정의한 SVI는 교전 전 공개정보로 어느 정도 예측할 수 있으며, 시작 승률과 경기 시각을 넘어 어떤 관측 정보가 추가 예측력을 제공하는가?
  \item[M-RQ4 (예측의 의미와 적용 범위)] 구성한 교전 결과 예측은 초기 경기 우세에 얼마나 의존하며, 균형 상태·경기 시간·관측 조건·패치와 지역이 달라질 때 어떤 범위까지 유지되는가?
\end{description}

M-RQ1은 제 \ref{ch:data} 장과 \ref{ch:engagement} 장이, M-RQ2는 제 \ref{ch:value} 장이, M-RQ3은 제 \ref{ch:prediction} 장이, M-RQ4는 제 \ref{ch:validation} 장이 답하며, 각 답의 적용 조건은 해당 결과의 첫 설명에, 공통 해석 범위는 제 \ref{sec:disc-scope} 절에 둔다.
